%CCP 2010, Chimie 2, CCP

\setcounter{equation}{0}
\setcounter{Q}{0}

Le plus souvent les procédés d'élaboration pyrométallurgiques conduisent à des métaux contenant des impuretés avec lesquelles peuvent se former des solutions solides. Cette partie propose une première approche thermodynamique des propriétés des solutions solides.

On considère dans toute cette partie le cas d'une mole d'un alliage formé par la solution solide entre $n_A$ mole de À et $n_B$ mole de B, soit $n_A + n_B = 1$.

\paragraph*{Enthalpie libre d'un mélange binaire}
\Q  L'enthalpie libre est-elle une grandeur intensive ou extensive ? Donner trois autres exemples de fonctions et/ou variables d'état du même type.

\Q Rappeler la définition du potentiel chimique $\mu_A$ de l'espèce \ce{A} en fonction de l'enthalpie libre $G$ du système fermé constitué de $n_A$ mole de \ce{A} et $n_B$ mole de \ce{B}.

Comment s'exprime $G$ en fonction des potentiels chimiques des $2$ espèces ?

Donner l'expression de l'identité thermodynamique donnant la différentielle de $G$, notée $dG$, pour ce système binaire fermé en fonction des variables $P$, $T$ et des quantités de matière $n_A$ et $n_B$.

\Q \label{Q:01-03} Dans le cas du système fermé décrit ci-dessus et à $T$ et $P$ constantes, démontrer que
\begin{align*}
\mu_A &= G + (1-n_A) \frac{dG}{dn_A} \\
\mu_B &= G - n_A \frac{dG}{dn_A}
\end{align*}

%\Q Reproduire sur votre feuille l'allure de G = f(n4) (à T et P constants) donnée ci-dessous.
\Q \label{Q:01-04} Proposer une méthode de détermination graphique de $\mu_A$ et $\mu_B$ : l'illustrer sur la courbe $G = f(n_A)$ pour un paramètre de composition $n_A$ quelconque (point $M$).

\paragraph*{Grandeur de mélange et solutions solides idéales}

On appelle grandeur de mélange $\Delta X_m$, la variation de la grandeur $X$ lors de la formation isobare et isotherme d'une solution à partir des composés pris purs dans leur état le plus stable. Dans le cas d'une mole de solution solide où \ce{A} et \ce{B} sont solides dans leur état pur standard, $\Delta X_m$ s'identifie à la variation de $X$ pour la réaction suivante :

\begin{align*}
\ce{n_A A(s) + n_B B(s) = (A_{n_A}B_{n_B})(s)}
\end{align*}

\Q \label{Q:01-05} Soit $\mu_A^\star$ et $\mu_B^\star$ les potentiels chimiques des corps purs. Exprimer $\Delta G_m$ en fonction de l'enthalpie libre de la solution solide $G$, $\mu_A^\star$, $\mu_B^\star$ et $n_A$.

\Q Déterminer l'expression de $\Delta G_m$ en fonction de la température, de $n_A$ et des activités $a_A$ et $a_B$ des espèces \ce{A} et \ce{B} dans la solution.

\Q Que devient l'expression précédente de $\Delta G_m$ dans le cas d'une solution idéale ?

Pour $n_A = 0$ et $n_A = 1$, calculer $\Delta G_m$ et déterminer les pentes des tangentes à la courbe $\Delta G_m = f(n_A)$.

Quelle est la valeur de $n_A$ pour laquelle $\Delta G_m$ est extremum ?

\Q En déduire l'allure de la courbe $\Delta G_m = f( n_A )$ dans le cas d'une solution idéale et justifier pourquoi il est difficile d'obtenir des matériaux de très haute pureté.

En vous aidant des questions \ref{Q:01-05} et \ref{Q:01-03}, préciser les grandeurs qui peuvent être déterminées à partir de la courbe $\Delta G_m  = f ( n_A )$ par une exploitation graphique analogue à celle développée à la question \ref{Q:01-04}.

L'enthalpie $\Delta H_m$ et l'entropie $\Delta S_m$ de mélange sont reliées à l'enthalpie libre de mélange par la relation :

\begin{align*}
\Delta G_m = \Delta H_m - T \Delta S_m
\end{align*}

\Q Déterminer les expressions de $\Delta S_m$ et $\Delta H_m$ dans le cas d'une solution idéale.

Quelle est la valeur de $n_A$ pour laquelle $\Delta S_m$ est extremum ? S'agit-il d'un minimum ou d'un
maximum ? Commenter.

\Q On considère une solution solide où \ce{A} et \ce{B} pris purs cristallisent dans le système cubique avec Z atomes par maille : leurs paramètres de maille sont respectivement $a_A$ et $a_B$ à l'état
pur.
\begin{enumerate}
\item Calculer le volume de mélange $\Delta V_m$ associé à la formation de la solution solide supposée idéale, sachant que : 

\begin{align*}
	\Delta V_m = \left( \frac{\partial \Delta G_m}{\partial P} \right)_{T, n_A}
\end{align*}

\item En déduire une relation entre le paramètre de maille $a$ de la solution solide, $a_A$ et $a_B$.

Sachant que $a_B = (1+\epsilon) \cdot a_A$ avec $\epsilon \ll 1$, on admettra que la relation précédente permet de démontrer la loi de Végard qui stipule une évolution linéaire du paramètre de maille $a(n_A)$ de la solution solide entre $a_A$ et $a_B$ en fonction du paramètre de composition de l'alliage $n_A$.

\end{enumerate}

Le système nickel - cuivre est un exemple quasi-parfait de solution idéale. Le nickel cristallise dans le système cubique en formant un empilement C.F.C. (cubique à faces centrées).

\Q Représenter la maille du nickel pur. Donner le nombre d'atomes par maille $Z$ et la coordinence $N$ des atomes de nickel.

Les rayons du nickel et du cuivre sont respectivement $r_{Ni} = 125 ~ \si{\pico\meter}$ et $r_{Cu} = 128 ~ \si{\pico\meter}$. Est-ce que ces deux métaux conduisent à des solutions solides de substitution ou d'insertion ? Justifier.

\Q Quel est selon vous le type structural du cuivre pur ? Justifier.

\Q On considère une mole de solution solide constituée de $n_{Ni}$ mole de \ce{Ni} et $n_{Cu}$ mole de \ce{Cu}, et on note $a$ le paramètre de maille de la solution solide. En attribuant un rayon moyen aux atomes de la solution solide, qui tient compte de la composition de l'alliage, déterminer la relation entre $a$, $n_{Ni}$ et les paramètres de maille $a_{Ni}$ et $a_{Cu}$ des corps purs \ce{Ni} et \ce{Cu}. Ce résultat constitue la loi de Végard.
